%
% Cover letter for AIP Advances revised submission (MS #ADV26-AR-04146)
% (see AIP_Style_4thed.pdf, Sec. II.D, on including a cover letter)
%
\documentclass[11pt]{letter}

\usepackage[T1]{fontenc}
\usepackage{lmodern}
\usepackage[margin=1in,top=0.75in,bottom=0.75in]{geometry}

% Sender address block (corresponding author)
\address{
Hideo Takahashi \\
Managed \& Platform Services Business Division, Hitachi, Ltd. \\
292, Yoshida, Totsuka-ku, Yokohama, 244-0817, Japan \\
hideo.takahashi.qa@hitachi.com
}

\signature{Hideo Takahashi}

\begin{document}

\begin{letter}{Dr. Tiange Bi \\ Associate Editor \\ AIP Advances \\ AIP Publishing}

\opening{Dear Dr. Bi,}

Thank you for the opportunity to revise our manuscript, MS \#ADV26-AR-04146,
``Grid-Based Quantum Circuit Simulation of the Hydrogen Atom: Coulomb
Singularity Regularization and Parameter-Selection Guidelines,'' by Hideo
Takahashi, Tatsuya Tomaru, Toshiyuki Hirano, Saisei Tahara, and Fumitoshi Sato,
for consideration in AIP Advances. We thank Reviewer \#1 for the careful
reading and constructive comments. We have revised the manuscript accordingly
and provide a point-by-point response below, with page numbers referring to
the revised manuscript.

\begin{enumerate}

\item \textit{The writing of the abstract could be further improved. For
example, the placement of (1), (2), and (3) in the last sentence of the
abstract is somewhat awkward. If these are intended as parallel items, each
should fit the same grammatical structure; otherwise, the numbers could simply
be omitted. The earlier list of contributions, (1)--(3), would also benefit
from a brief introductory sentence, as the enumeration currently begins rather
abruptly.}

We have removed the (1)/(2)/(3) numbering throughout the abstract and instead
refer back to each contribution by its content (e.g., ``the circuit
implementation,'' ``the singularity treatment,'' ``the parameter-selection
procedure''), which avoids both the awkward parallel-structure issue in the
final sentence and the abrupt start of the enumeration. We also removed a
repeated occurrence of ``one- and two-dimensional hydrogen atom,'' replacing
the second instance with ``these systems.'' (Abstract, p.~1.)

\item \textit{I am not sure some wording about the sampling-theorem argument
in Section II D 2 is fully accurate. In particular, the Nyquist constraint may
not always be necessary for accurate Trotter simulation. One naive example is
when $H_1$ and $H_2$ commute, then time step can be anything (not restricted
by sampling constraints). Although Coulomb Hamiltonian terms certainly do not
commute, it would still be beneficial to be careful in wording/argument when
justifying this sampling constraint as mandatory.}

We thank the reviewer for this comment. We agree that the sampling constraint
is not required when $H_1$ and $H_2$ commute, since the decomposition is then
exact and larger phase changes per step are allowed. The constraint is thus
specific to the non-commuting case, which applies to the Coulomb Hamiltonians
considered in this work. We have revised Sec.~II~D~2 to state this
explicitly, citing Kosloff [J.~Phys.~Chem.\ \textbf{92}, 2087 (1988)], who
pointed out that such a restriction on the time step applies to short-time
propagators that factorize the evolution operator into non-commuting parts.
We have also clarified that the phase condition is imposed on each of the
decomposed Hamiltonians. The conditions themselves [Eqs.~(54)--(55)] are
unchanged, so all derived time-step limits and numerical results remain the
same.

\smallskip
\noindent\textbf{Submitted version:}
\begin{quote}
\itshape
``The first condition that the time step $\delta t$ must satisfy is the
avoidance of aliasing errors derived from the sampling theorem for the
Discrete Fourier Transform. Considering one step of time evolution via the
Suzuki--Trotter decomposition as a sampling process, the amount of phase
change by the Hamiltonian must not exceed the phase change corresponding to
the Nyquist frequency; that is, the maximum phase change per step must be
less than a half-period ($\pi$). This requires the following inequalities for
the potential and kinetic energy terms, respectively: [Eqs.~(54)--(55)]''
\end{quote}

\noindent\textbf{Revised version:}
\begin{quote}
\itshape
``The first condition that the time step $\delta t$ must satisfy is the
avoidance of aliasing errors derived from the sampling theorem for the
Discrete Fourier Transform. Considering one step of time evolution via the
Suzuki--Trotter decomposition as a sampling process, the amount of phase
change by each of the decomposed Hamiltonians must not exceed the phase change
corresponding to the Nyquist frequency; that is, the maximum phase change per
step must be less than a half-period ($\pi$). This constraint is specific to
the non-commuting case: if $H_p$ and $H_k$ commuted, the decomposition would
be exact and larger phase changes per step would be allowed [Kosloff 1988].
Since $H_p$ and $H_k$ do not commute in general, as in the case of the
hydrogen atom, we require the following inequalities for the potential and
kinetic energy terms, respectively: [Eqs.~(54)--(55), unchanged]''
\end{quote}
(Sec.~II~D~2, p.~12.)

\item \textit{Given the relevance of high-order Trotter error for Coulomb
systems and the theoretical understanding of improved convergence and angular
momentum for the hydrogen atom, arXiv:2604.07704 should also be cited.}

We thank the reviewer for pointing us to this work by Fang and Wu. We have
cited it in the Introduction, following the discussion of the $n^{-1/4}$
scaling for Coulomb systems, and briefly noted its result that the
convergence order expected for bounded operators is recovered for
hydrogen-atom initial states with sufficiently high orbital angular momentum.

\smallskip
\noindent\textbf{Submitted version} (Sec.~I, p.~3):
\begin{quote}
\itshape
``$\dots$ the error increase relative to the particle number is polynomial.
In our study, using numerical evaluations based on an emulator and classical
programs, we present an integrated approach for a single-electron model:
$\dots$''
\end{quote}

\noindent\textbf{Revised version:}
\begin{quote}
\itshape
``$\dots$ the error increase relative to the particle number is polynomial. A
subsequent study showed that the convergence order expected for bounded
operators is recovered for hydrogen-atom initial states with sufficiently high
angular momentum $\ell$: $\ell \ge 2$ for the first-order and $\ell \ge 4$ for
the second-order formula [Fang and Wu, arXiv:2604.07704]. In our study, using
numerical evaluations based on an emulator and classical programs, we present
an integrated approach for a single-electron model: $\dots$''
\end{quote}

\item \textit{On p.~22, I am not sure the wording that ``Because states with
non-zero angular momentum are smooth at the origin and lack a cusp, the
divergence of the higher moments is avoided'' for $\psi^{2D}_{1,1}$ is fully
accurate, as $\|V^3\psi^{2D}_{1,1}\|_{L^2}$ (and $\|V^2\psi^{2D}_{1,1}\|_{L^2}$
both) diverge, making the corresponding term in Eq.~(73) divergent. Please
revise the wording in this part to enhance accuracy.}

We thank the reviewer for this important correction and have substantially
revised this passage. We now state precisely that states with non-zero angular
momentum vanish at the origin as $\psi\sim r^{|m|}$ rather than being smooth,
so that the singularity of the terms in the state-dependent error bound is
weakened but not removed. In agreement with the reviewer's calculation, we now
state that in the continuum limit $\delta r\to0$ these terms diverge for
$\psi^{2D}_{1,1}$: $\|V^3\psi\|$ remains finite only for $|m|\ge3$, and the
terms involving the kinetic energy require even larger $|m|$. On our grid,
every term is finite, because the potential is capped at $|V(0)|=4/\delta r$
and the kinetic energy is bounded by the maximum wavenumber of the grid.
However, since this also holds for $m=0$ states, finiteness alone does not
explain the difference between $m=0$ and $m=1$. We therefore base the
explanation on how fast the terms grow as $\delta r$ decreases: a
power-counting estimate based on $\psi\sim r^{|m|}$ near the origin indicates
that each term grows by one power of $1/\delta r$ less for $|m|=1$ than for
$m=0$ (e.g., $\|V^3\psi\|\propto\delta r^{-1}$ versus $\delta r^{-2}$), so that
at the resolution used here the second-order bound is considerably smaller for
$\psi^{2D}_{1,1}$, consistent with the observed advantage. We also state
explicitly that this advantage relies on the finite resolution and is not
guaranteed to persist in the continuum limit for $|m|=1$.

\smallskip
\noindent\textbf{Submitted version:}
\begin{quote}
\itshape
``Because states with non-zero angular momentum are smooth at the origin and
lack a cusp, the divergence of the higher moments is avoided. Consequently,
the inherent error-reduction benefit of the higher-order decomposition
functions as expected, manifesting as the significant difference observed in
accuracy.''
\end{quote}

\noindent\textbf{Revised version:}
\begin{quote}
\itshape
``Because states with non-zero angular momentum vanish at the origin
($\psi\sim r^{|m|}$) rather than forming a cusp, the terms of Eq.~(73) are
less singular than for $m=0$ states; the singularity is weakened, however,
not removed. In the continuum limit, $\|H_2(g)^3\phi\|$ remains finite only
for $|m|\ge3$, and the terms involving the kinetic energy require even larger
$|m|$, so that for $\psi^{2D}_{1,1}$ the bound diverges as $\delta r\to0$. On
the grid, by contrast, every term is finite, because the potential is capped
at $|V(0)|=4/\delta r$ and the kinetic energy is bounded by the maximum
wavenumber of the grid. Since this holds equally for $m=0$ states, what
distinguishes the two cases is how fast the terms grow as $\delta r$
decreases. A power-counting estimate based on the behavior $\psi\sim r^{|m|}$
near the origin indicates that each term grows by one power of $1/\delta r$
less for $|m|=1$ than for $m=0$; for example, $\|H_2(g)^3\phi\|$ grows as
$\delta r^{-1}$ for $\psi^{2D}_{1,1}$ but as $\delta r^{-2}$ for
$\psi^{2D}_{0,0}$, and the same one-power reduction applies to the mixed term
$\|H_2(g)H_1(g)^2\phi\|$ and the kinetic term $\|H_1(g)^3\phi\|$. At the
resolution used here, the second-order bound is therefore expected to be
considerably smaller for $\psi^{2D}_{1,1}$, which is consistent with the
observed advantage of the second-order decomposition. Because this advantage
relies on the finite resolution, it is not guaranteed to persist in the
continuum limit for $|m|=1$.''
\end{quote}
(Sec.~IV, p.~21.)

In addition, since this discussion relies on the state-dependent Trotter error
bound of Burgarth et al.\ already used in our manuscript, we also clarified
our identification of their operators $H_1$ and $H_2$ with the kinetic and
potential energy terms, respectively, which was previously left implicit:

\smallskip
\noindent\textbf{Submitted version} (Sec.~II~D~3, p.~13):
\begin{quote}
\itshape
``$\dots$ this study introduces the state-dependent Trotter error bound
proposed by Burgarth et al., which assumes a specific target wave function.''
\end{quote}

\noindent\textbf{Revised version:}
\begin{quote}
\itshape
``$\dots$ this study introduces the state-dependent Trotter error bound
proposed by Burgarth et al., which assumes a specific target wave function.
Following their convention we identify $H_1=H_{\mathrm{ek}}$ and
$H_2=H_{\mathrm{ep}}$.''
\end{quote}

\end{enumerate}

We believe these revisions fully address the reviewer's comments and have
improved the accuracy and clarity of the manuscript. We hope the revised
manuscript is now suitable for publication in AIP Advances.

Thank you for your consideration.

\closing{Sincerely,}

\end{letter}
\end{document}
